% ============================================================
% 摘要
% ============================================================

\begin{cabstract}
% 在此撰写中文摘要内容
% 摘要应概括论文的主要内容，包括研究目的、方法、结果和结论
% 一般为 500-1000 字

网络与通信技术的迅速发展给人们的工作和生活方式带来了巨大的改变。在更大、更快、更安全、更方便且融合有线传输和无线传输于一体的下一代互联网NGI（Next Generation Internet）中，随时随地享受高质量网络服务已经成为人们的迫切要求，这在客观上要求NGI提供可靠的QoS（Quality of Service）保证，并且可以在通信开始和运行期间为用户提供总最佳连接ABC（Always Best Connected），进而使用户总是以最优方式接入网络并享受服务。

\end{cabstract}

\ckeywords{下一代互联网；总最佳连接；切换；服务质量；效用；遗传算法
}

\begin{eabstract}
% Write your English abstract here
% The abstract should summarize the main content of the thesis

The rapid development of Internet and communication technology has brought tremendous changes to the way people work and live. In a larger, faster, more secure, more conventient NGI(Next Generation Internet) that integrates wired and wireless access technologies, enjoying high-quality network service anywhere and at any time has become people’s necessary requirement, which indicates that NGI must provide QoS(Quality of Service) guarantee and support ABC(Always Best Connected). Therefore, people can choose the “best” way to access to network and enjoy the service.

\end{eabstract}

\ekeywords{next generation internet; always best connected; handoff; quality of service; utility; genetic algorithm}
